On considère la fonction $g$ définie sur l'intervalle $[0~;~1]$ par 

\[g(x) = 2x - x^{2}.\]

\smallskip

\begin{enumerate}
\item Montrer que la fonction $g$ est strictement croissante sur l'intervalle $[0~;~1]$ et préciser les valeurs de $g(0)$ et de $g(1)$.
\end{enumerate}

On considère la suite $\left(u_{n}\right)$ définie par $\left\{\begin{array}{l c l}u_{0}&=&\dfrac{1}{2}\\ u_{n+1}&=&g\left(u_{n}\right)\end{array}\right.$ pour tout entier naturel $n$.

\begin{enumerate}
\setcounter{enumi}{1}
\item Calculer $u_{1}$ et $u_{2}$.

\item Démontrer par récurrence que, pour tout entier naturel $n$, on a : $0 < u_{n} < u_{n+1}< 1$.
\item En déduire que la suite $\left(u_{n}\right)$ est convergente.
\item Déterminer la limite $\ell$ de la suite $\left(u_{n}\right)$.
\end{enumerate}

On considère la suite $\left(v_{n}\right)$ définie pour tout entier naturel $n$ par $v_{n}=\ln \left(1 - u_{n}\right)$.

\begin{enumerate}
\setcounter{enumi}{5}
\item Démontrer que la suite $\left(v_{n}\right)$ est une suite géométrique de raison 2 et préciser son premier terme.
\item En déduire une expression de $v_{n}$ en fonction de $n$.
\item En déduire une expression de $u_{n}$ en fonction de $n$ et retrouver la limite déterminée à la question 5.
\item Recopier et compléter le script Python ci-dessous afin que celui-ci renvoie le rang $n$ à partir duquel la suite dépasse 0,95.
\end{enumerate}

\begin{center}
\begin{tabular}{|>{\texttt} l|}\hline

\textbf{def }seuil():\\
\qquad n=0\\
\qquad u=0.5\\
\qquad \textbf{while} u < 0.95 :\\
\qquad \qquad n=...\\
\qquad \qquad u=...\\
\qquad \textbf{return} n\\ \hline
\end{tabular}

\end{center}
